\section{A positive bound for the final two layers}\label{sec:positiva}

We first apply double counting in one direction only. The deterministic bound on the number of reverse operations turns a lower bound on the
mean number of forward repairs into the lower bound on $f_h$ in Theorem~\ref{thm:positiva}. The early cut-off supplies the repairs
under the reference process. It remains to determine what proportion of the layer retains this bound under the uniform law.

Let $h\in\{1,2\}$ and let $\mathcal B_n\subseteq\Om_h(n)$ be the event that some empty crossing has fewer than $a_0n/2$ legal one-move
repairs (including the $G$ filter).

\emph{Under the process.} Let $\omega$ be a successful trajectory on which~\eqref{eq:corte} holds. For every crossing that is empty at
the end, Lemma~\ref{lem:perdida} with $\delta=\delta_0$ and~\eqref{eq:constantes} give \[
 A^{rc}_{\mathrm{final}}\ge a_0n-4\delta_0n-4\ge a_0n/2
\] for large $n$, since $a_0>8\delta_0$. The bound~\eqref{eq:corte} holds simultaneously for all roots before we select those that
remain empty, so there is no survivorship bias. By Lemma~\ref{lem:corte}, \[
 \QQ_h(S\cap\{\text{the final placement lies in }\mathcal B_n\})\le e^{-cn}.
\] Failure of the full process need not be rare.

\emph{Transfer.} By Theorem~\ref{thm:comparacion}, $\QQ_h(S)\ge e^{-d_n}$, so $\nu'_h(\mathcal B_n)\le e^{-cn+d_n}\le e^{-cn/2}$ for
large $n$. Lemma~\ref{lem:binaria} and $\KL{\mu_h}{\nu'_h}\le d_n=o(n)$ give \[
 \mu_h(\mathcal B_n)\le\frac{d_n+\log2}{cn-d_n}\longrightarrow0 .
\] This gives convergence to zero under $\mu_h$ without an exponential tail bound.

\emph{Double counting.} Each placement in $\Om_h$ has exactly $h^2$ empty crossings, and $d_{h,1}$ sums the repairs over all of them.
On $\mathcal B_n^c$, $d_{h,1}\ge h^2a_0n/2$. By Lemma~\ref{lem:doblecuenta}, \[
 f_h\ge\frac{\E_{\mu_h}d_{h,1}}{h^2(n-h)}\ge\frac{a_0n}{2(n-h)}\,\mu_h(\mathcal B_n^c),
 \qquad\text{hence}\qquad \liminf_{n\to\infty}f_h(n)\ge\frac{a_0}2 .
\] Normalizing by the $h^2$ crossings avoids losing a factor of $h^2$. \qed

\begin{observacion} The result is asymptotic and provides no useful constant for small boards. We do not replace $a_0$ with the
numerical mean of the profile. The proof uses the reference process as defined, without conditioning its transitions during the
analysis or taking any product of vacancies under $\mu_h$. Here the number of queens placed after the cut-off is $m_{\delta_0}\asymp
n$. \end{observacion}
